\chapter{Conclusion and Future Scope}\label{ch:conc}

\section{Conclusion}
This project built PaperMind, an end-to-end system that learns the pattern of an examination from its own past papers and
uses that pattern to generate new, verified practice papers and to coach individual students. The NLP side ingests
digital, scanned and unfamiliar papers; segments and joins GATE questions with their official keys; labels every
question with a syllabus topic through an LLM teacher distilled into a lightweight classifier; assigns Bloom's levels;
and reports subject weightage, recency-weighted hot topics, trends, question-type mix and concept recurrence. On a
documented corpus of 31 papers and 1,971 questions, the distilled classifier reached 77.7\% topic accuracy across 58
topics and 88.5\% subject accuracy on unseen 2026 GATE CS papers, outperforming a zero-shot embedding baseline while
running without any neural network.

The GenAI side turns a one-line request into a structured specification, infers a blueprint from history, samples
topics in proportion to their importance, retrieves past questions as exemplars and lets Gemini write the paper. Every
paper is validated for structure, marks and novelty, and every objective answer key is re-solved independently and
replaced if it fails---a safeguard motivated by a real incorrect key found during testing. A complete GATE CS mock is
produced in about three minutes, matches the historical subject distribution with $r=0.81$, and contains no question
closer than 0.61 to any past question. A tool-using assistant grounds its advice in the corpus, quizzes students and
prioritises topics by examination weight and measured weakness. The whole system is deployed as a web application at no
cost.

Beyond the artefact, the project yielded practical findings: text-layer extraction destroys mathematical notation;
multimodal LLMs refuse to transcribe widely published papers; public mirrors can mislabel files (an Electronics key filed
as Computer Science); LLM-written answer keys must be checked; and free LLM tiers impose request budgets that shape the
architecture.

\section{Future Work}
\begin{itemize}
  \item \textbf{More subjects and exams.} Extend the corpus to other GATE papers (EC, EE, ME, MA) and to university
        papers of each course, which the pattern-agnostic ingestion already supports.
  \item \textbf{Human-validated labels.} Annotate a sample of questions to measure true labelling accuracy and to
        fine-tune the student model.
  \item \textbf{Stronger verification.} Execute code-based questions, solve numerical questions symbolically or with
        program-aided reasoning, and use several solver samples to estimate confidence.
  \item \textbf{Figures and diagrams.} Extract diagrams from past papers and generate questions with figures (circuits,
        automata, trees) using a drawing language such as TikZ or Graphviz.
  \item \textbf{Difficulty calibration.} Estimate question difficulty from student responses (item response theory)
        and generate papers at a target difficulty.
  \item \textbf{Persistent, multi-user deployment.} Add accounts and a managed database so that progress survives
        restarts, and adaptive spaced-repetition schedules in the assistant.
  \item \textbf{Mobile and offline use.} Package practice sets for offline use in low-connectivity regions.
\end{itemize}

\section{Chapter Summary}
PaperMind meets its objectives: it analyses past papers objectively, generates pattern-faithful and verified practice
papers on demand, and personalises preparation through an assistant, all deployable for free. The limitations
identified---verification that is strong but not absolute, the absence of figures, and labels measured against an LLM
teacher---define a clear path for future work.
